\documentclass[11pt, a4paper]{article}
\usepackage[margin=1in, top=1.2in, bottom=1.2in]{geometry}
% ---- Font & Micro-Typography ----
\usepackage[T1]{fontenc}
\usepackage{lmodern}
\usepackage{microtype}
\usepackage{amsmath}
% ---- Formatting & Aesthetics ----
\usepackage{graphicx}
% ---- Bibliography ----
\usepackage[backend=biber, style=numeric-comp, sorting=nyt, giveninits=true]{biblatex}


% ====================================================================================
\begin{document}

\begin{center}
  \LARGE\textbf{Daily Rant on Business and Entrepreneurship}\\
  Tadouanla Guetchuin Billy\\
  28177
\end{center}

\section{Begin}
\begin{quote}
  \textit {Today it is more about Sustainable businesses rather than simple profit making
    mentality.}\\

  \textit {The lab for business Entrepreneurship is the actual world ( \textbf{Real world})
    just like the lab for this study program is the \textbf{Spoy canal}.}\\

  \textit{We are mainly going to look into what makes a business a business as well as the
    sustainability aspect of businesses.}
\end{quote}

\section{What makes a Business a Business}
When we are talking about a business we are talking about dealing with a stream of income in
a way.

\begin{quote}
  \itshape
  Not having a boss is one of the main driving factors for people starting a
  business
\end{quote}

\begin{quote}
  \itshape
  A business cannot grow forever; every market has a ceiling.
\end{quote}

\section{Profit Maximization: The Main Goal of Any Business}
Profit is the difference between revenue and cost, both of which depend on
the level of output:

\begin{center}
  \text{Profit} = \text{Total Revenue} - \text{Total Cost}
\end{center}

\subsection{Marginal Revenue and Marginal Cost}
\begin{itemize}
  \item \textbf{Marginal revenue (MR):} the addition to total revenue from
        producing one more unit of output. When a firm is a price taker,
        marginal revenue equals the price.
  \item \textbf{Marginal cost (MC):} the addition to total cost from
        producing one more unit of output.
  \item Profit is maximized at the output where MR = MC.
  \item Example: a cigarette-producing machine must be fully used to achieve
        a low unit cost.
        \begin{itemize}
          \item Spreading the fixed cost of the machine over more units
                lowers the cost per unit.
        \end{itemize}
\end{itemize}

\subsection{Market Share and Pricing}
\begin{itemize}
  \item Gaining market share is closely tied to offering cheaper prices.
  \item Scarcity pushes prices higher.
  \item Natural monopolies are rare: there is almost always some
        competition or substitute.
\end{itemize}

\subsection{Growth Limits and Diversification}
\begin{itemize}
  \item A business cannot grow forever. There is always a growth ceiling,
        and growth eventually plateaus.
  \item Staying with a single product inevitably leads to a plateau.
  \item Firms can keep growing by entering other markets, developing new
        products, or buying other businesses.
\end{itemize}


\begin{figure}[h!]
  \includegraphics[scale=0.4]{./figures/MR_MC.png}
\end{figure}


\noindent Think of including the graph of marginal cost and marginal Revenue. Marginal revenue
tends to be higher at the beginning and drops over time. The more you try to sell the less you
can as for the additional unit. When marginal cost equals marginal revenue you have to stop.
You now run the risk of increasing your marginal cost and dropping the marginal revenue even
more.

\begin{quote}
  \textit{You know breaking things down also has to do with the dumbing things down not just
    physically sectioning things if it happens that you can't follow or understand, dumb it
    down for you }\\

  \textit{When taking notes make sure to have everything that are related stay together rather
    than having multiple points. You could have them in paragraph form you know that will make
    the after reading more understandable.}
\end{quote}
%------------------- Work on this section -----------------------
\textbf{share price?}

\section{Profit maximization in the real world}
Profit Maximization very often does not guide business decision making. Some people are pretty okay with the current state of their business and are not willing to take the decision of
venturing into something new or taking risks.

Sales maximization is aimed at attempting to maximize the volume of sales rather than the
return gained from them. This could be done in order to establish a monopoly on the market.\\

\noindent\textbf{Share Price Maximization:} Pursuing policies aimed at increasing the share
price on the market.\\
\textbf{Profit Satisficing} Generating sufficient profits to satisfy shareholders but
maximizing\\
at a certain level business owners are satisfied with the whole idea of making money.\\
Learn more about corporations.
the rewards to the managers/board, and avoiding attention from rivals or regulatory authorities.

\begin{itemize}
  \item Convincing the Capital market of the share value of the product.
  \item Infinite growth is impossible.
  \item less than $0.05\%$
\end{itemize}

%------------Looking into this as well ---------------------
\begin{itemize}
  \item How many report  \textbf{more than 500m} in global revenue? 1,800 = less than 0.5\%
  \item How many are publicly traded corporations? 12,000 = 0.3\%
  \item how many report less than  \textbf{1m} in global revenue? 3.3 million = 87\%
  \item How many do some kind of manufacturing? roughly 630,000 = 16.6\%

  \item which business is the largest in Germany? \textbf{VW} 330 billion in revenue, 680,000
        employees.
  \item Which business is the largest privately owned in Germany? \textbf{SCHWARZ} group:
        190 billion in revenue, 600,000 employees\\
\end{itemize}

\noindent Boehringer family \\
Albrecht Aldi

\section{What is a Business Model}
\textbf{Definition:} A business model is the Entrepreneur's logic that will allow a venture to:
\begin{itemize}
  \item Capture the Market opportunity;
  \item Mitigate risks;
  \item identify the required resource set; and; and
  \item Create value for investors and founders.
\end{itemize}

\noindent\textbf{Describing The WHY:}
\begin{itemize}
  \item The business \textbf{idea bridges idea and action}.
        It answers the question of  \textbf{why} a venture will be viable or valuable.
  \item it is not a strategy and not a plan.
  \item It evolve overtime.
\end{itemize}

\subsection{Defining the Business Model}
Business models are formed through a process of addressing a series of questions
\begin{itemize}
  \item In which market do we operate?
  \item Which \textbf{product or service} do we offer
  \item Which part of the value created there will actually be created by us.
  \item How will we cooperate with other businesses in creating our value?
  \item Which resources do we need \textbf{just once}?
  \item Which resources do we need \textbf{Everyday}?
  \item How will our business evolve over time?
\end{itemize}

\begin{figure}
  \center
  \includegraphics[scale=0.4]{./figures/BM.png}
  \caption{The 6 various business models shown on the Hexagon above containing a series of
    questions to be answered in order to develop a proper business model}
\end{figure}

\end{document}
